% Ejemplo numerico completo de proyeccion ICOCIV: Vias urbanas, h=18.
% Documento autocontenido: compila con pdflatex + bibtex + pdflatex x2.
% Todos los valores provienen de una ejecucion real de la aplicacion.
\documentclass[11pt,letterpaper]{article}
\usepackage[utf8]{inputenc}
\usepackage[T1]{fontenc}
\usepackage[spanish,es-noquoting]{babel}
\usepackage[margin=2.4cm]{geometry}
\usepackage{amsmath}
\usepackage{booktabs}
\usepackage{longtable}
\usepackage{array}
\usepackage{float}
\usepackage{setspace}
\usepackage[hidelinks]{hyperref}

\title{Ejemplo numérico de proyección ICOCIV\\
\large Serie \textit{Vías urbanas}, horizonte solicitado de 18 meses}
\author{Proyecto ICOCIV / ICCP--ICOCIV}
\date{\today}

\begin{document}
\maketitle

\begin{abstract}
\noindent
Este documento reproduce, paso a paso y con cálculos verificables, cómo la
aplicación obtiene el valor proyectado del índice ICOCIV para la serie
\textit{Vías urbanas} con un horizonte solicitado de 18 meses. Todos los
valores provienen de una ejecución real sobre el anexo
\texttt{anex-ICOCIV-may2026.xlsb}; no se emplean datos ilustrativos. El
resultado final es un índice de \textbf{156{,}9057} para
\textbf{noviembre de 2027}, entregado como \textbf{escenario de alta
incertidumbre}, con un máximo recomendado de 12 meses.
\end{abstract}

\tableofcontents
\newpage

\section{Introducción del ejemplo}

La aplicación toma una ruta jerárquica del anexo ICOCIV, construye la serie
mensual asociada, valida su calidad, compara modelos mediante backtesting
walk-forward, decide hasta qué horizonte la evidencia respalda una proyección y
finalmente calcula la trayectoria. Este ejemplo recorre esas etapas con una
serie concreta.

La aplicación apoya el criterio técnico del profesional responsable; no lo
sustituye ni reemplaza la producción estadística oficial del
DANE \cite{dane_metodologia_icociv}.

\section{Selección de la serie ICOCIV}

La ruta jerárquica determina la tabla de origen y la fila fuente de la que se
extraen los valores mensuales.

\begin{table}[H]
\centering
\caption{Ruta ICOCIV utilizada en el ejemplo}
\begin{tabular}{@{}p{4.2cm}p{8.6cm}p{2.2cm}@{}}
\toprule
\textbf{Nivel} & \textbf{Valor seleccionado} & \textbf{Código} \\
\midrule
Grupo de obra & Carreteras, calles, vías férreas y pistas de aterrizaje, puentes, carreteras elevadas y túneles & 530201 \\
Subclase CPC & Carreteras (excepto carreteras elevadas); calles & 53211 \\
Tipología de obra & \textbf{Vías urbanas} & 53211\_07 \\
Tabla de origen & A16.2 Índice por tipologías de obra & \texttt{T\_16\_2} \\
\bottomrule
\end{tabular}
\end{table}

\begin{table}[H]
\centering
\caption{Identificación de la serie}
\begin{tabular}{@{}ll@{}}
\toprule
\textbf{Campo} & \textbf{Valor} \\
\midrule
Archivo fuente & \texttt{anex-ICOCIV-may2026.xlsb} \\
Identificador de fila fuente & 6 \\
Frecuencia & Mensual \\
Periodo inicial & 2021-01 \\
Periodo final & 2026-05 \\
Número de observaciones & 65 \\
Primer valor observado \(y_1\) & 100{,}801 \\
Último valor observado \(y_T\) & 144{,}590 \\
\bottomrule
\end{tabular}
\end{table}

\section{Serie histórica utilizada}

La serie tiene 65 observaciones mensuales continuas. Se muestran los primeros y
últimos seis periodos; la serie completa se adjunta en el archivo
\texttt{docs/ejemplo\_vias\_urbanas\_proyeccion.csv}.

\begin{table}[H]
\centering
\caption{Extremos de la serie histórica (índice y variación mensual)}
\begin{tabular}{@{}lrr@{}}
\toprule
\textbf{Periodo} & \textbf{Índice} & \textbf{Variación mensual (\%)} \\
\midrule
2021-01 & 100{,}801 & --- \\
2021-02 & 100{,}955 & 0{,}153 \\
2021-03 & 101{,}574 & 0{,}613 \\
\multicolumn{3}{c}{\dots} \\
2026-03 & 143{,}840 & --- \\
2026-04 & 144{,}310 & 0{,}327 \\
2026-05 & \textbf{144{,}590} & 0{,}194 \\
\bottomrule
\end{tabular}
\end{table}

\section{Preparación y validación de datos}

Antes de modelar, la aplicación normaliza los periodos a una secuencia mensual
ordenada y aplica un conjunto de validaciones. Un error crítico impide modelar;
una condición no crítica se registra como advertencia y permanece en la
trazabilidad.

\begin{longtable}{@{}p{3.4cm}p{5.6cm}p{2.4cm}p{4.0cm}@{}}
\caption{Validaciones aplicadas a la serie \textit{Vías urbanas}}\\
\toprule
\textbf{Validación} & \textbf{Criterio de la aplicación} & \textbf{Resultado} & \textbf{Observación} \\
\midrule
\endfirsthead
\toprule
\textbf{Validación} & \textbf{Criterio} & \textbf{Resultado} & \textbf{Observación} \\
\midrule
\endhead
Orden temporal & Periodos estrictamente crecientes & OK & Sin reordenamientos. \\
Continuidad mensual & Sin meses intermedios ausentes & OK & 65 meses consecutivos. \\
Valores faltantes & Conteo de nulos o no numéricos & 0 & No se imputa nada. \\
Valores duplicados & Periodos repetidos & 0 & --- \\
Valores no numéricos & Conversión numérica & 0 & --- \\
Valores no positivos & \(y_t>0\) & 0 & Requisito de las transformaciones logarítmicas. \\
Longitud mínima & \(\geq 8\) para modelar; \(\geq 24\) recomendado & OK (65) & Muy por encima del mínimo. \\
Atípicos robustos & \(|M_i|>3{,}5\) sobre variaciones & Detectados & Se reportan; no se eliminan \cite{nist_outliers}. \\
\bottomrule
\end{longtable}

La serie queda habilitada para modelación (\texttt{valida\_modelacion = True}).
La detección de atípicos marcó variaciones inusuales, lo que activó el
segundo nivel de modelos (Sección~\ref{sec:modelos}), pero no bloqueó el
análisis: los datos oficiales no se modifican.

\section{Análisis descriptivo inicial}

\begin{table}[H]
\centering
\caption{Estadísticos descriptivos de la serie}
\begin{tabular}{@{}lr@{}}
\toprule
\textbf{Estadístico} & \textbf{Valor} \\
\midrule
Mínimo & 100{,}801 \\
Máximo & 144{,}590 \\
Media & 121{,}658 \\
Mediana & 123{,}970 \\
Desviación estándar & 13{,}392 \\
Último valor observado & 144{,}590 \\
Variación mensual promedio & 0{,}569\% \\
Desviación de la variación mensual & 0{,}854\% \\
Variación acumulada del periodo & 43{,}441\% \\
\bottomrule
\end{tabular}
\end{table}

La variación acumulada se obtiene directamente de los extremos:
\begin{equation}
v_{\text{acum}}=\left(\frac{y_T}{y_1}-1\right)\times 100
=\left(\frac{144{,}590}{100{,}801}-1\right)\times 100 = 43{,}441\%.
\end{equation}

\paragraph{Interpretación.} La serie crece de forma sostenida durante todo el
periodo, con un avance medio cercano a \(0{,}57\%\) mensual y una dispersión
de las variaciones (\(0{,}85\%\)) mayor que la propia media. Es decir, hay
tendencia clara pero con ruido mes a mes apreciable: el nivel sube de manera
consistente, aunque el cambio de un mes concreto no es predecible con
precisión. Esa combinación anticipa que los métodos que extrapolan la tendencia
media de largo plazo tendrán ventaja sobre los que reaccionan al último
movimiento.

\section{Modelos candidatos evaluados}
\label{sec:modelos}

La aplicación activa modelos de forma progresiva. El nivel 1 siempre está
disponible; el nivel 2 se habilitó en este caso por la longitud de la serie y
por la presencia de atípicos relevantes. Los diez modelos evaluados fueron:

\begin{center}
\texttt{naive}, \texttt{drift}, \texttt{holt\_lineal}, \texttt{holt\_amortiguado},
\texttt{lineal}, \texttt{logaritmico}, \texttt{variacion\_lineal},
\texttt{exponencial\_log\_lineal}, \texttt{log\_variacion}, \texttt{huber}.
\end{center}

\begin{longtable}{@{}p{3.0cm}p{4.6cm}p{7.6cm}@{}}
\caption{Modelos candidatos, fórmula y rol en la aplicación}\\
\toprule
\textbf{Modelo} & \textbf{Fórmula} & \textbf{Rol y limitaciones} \\
\midrule
\endfirsthead
\toprule
\textbf{Modelo} & \textbf{Fórmula} & \textbf{Rol y limitaciones} \\
\midrule
\endhead
Naive & \(\hat y_{T+h}=y_T\) & Referencia mínima. No incorpora tendencia: subestima series con crecimiento sostenido \cite{hyndman2021fpp3}. \\
Drift & \(\hat y_{T+h}=y_T+h\,\dfrac{y_T-y_1}{T-1}\) & Benchmark con tendencia. Extiende la pendiente media histórica; es robusto pero ignora cambios de régimen. \\
Holt lineal & \(\hat y_{T+h}=\ell_T+h\,b_T\) & Nivel y tendencia suavizados; reacciona al comportamiento reciente \cite{holt1957}. \\
Holt amortiguado & \(\hat y_{T+h}=\ell_T+(\phi+\dots+\phi^{h})b_T\) & Modera la tendencia en horizontes largos \cite{gardner_mckenzie1985}. \\
Regresiones & OLS lineal, logarítmica, log-lineal, Huber & Referencias interpretables sobre el nivel; no son evidencia causal \cite{granger_newbold1974}. \\
Sobre variación & Modela \(v_t\) y reconstruye el nivel & Útil cuando la dinámica es más estable en cambios que en niveles; acumula error con el horizonte. \\
\bottomrule
\end{longtable}

\section{Backtesting walk-forward}

\subsection{Procedimiento}

Un \emph{origen} \(r\) es un punto de corte de la serie. Para cada origen, el
modelo se reestima usando \textbf{únicamente} \(\{y_1,\dots,y_r\}\) y pronostica
\(h\) meses adelante. El pronóstico se compara con el valor realmente observado:
\begin{equation}
e_{r,h}=y_{r+h}-\hat y_{r+h\mid r}.
\end{equation}
El origen avanza un mes y el proceso se repite. Como el ajuste de cada origen
solo ve el pasado, se evita por construcción la fuga de información
futura \cite{tashman2000}.

El entrenamiento inicial es \(\max(18;\,0{,}60n)\). Con \(n=65\):
\begin{equation}
\text{entrenamiento inicial}=\max(18;\ 0{,}60\times 65)=\max(18;\ 39)=39 .
\end{equation}
El número de orígenes disponibles para un horizonte \(h\) es
\(m(h)=n-39-h+1=27-h\); para \(h=1\) resultan \(26\) orígenes y para \(h=18\),
\(9\).

\subsection{Orígenes reales del modelo seleccionado (\(h=1\))}

\begin{table}[H]
\centering
\caption{Muestra de orígenes del backtesting de Drift con \(h=1\) (26 en total)}
\begin{tabular}{@{}llrrrr@{}}
\toprule
\textbf{Origen} & \textbf{Periodo evaluado} & \textbf{Obs. entren.} & \textbf{Observado} & \textbf{Pronosticado} & \textbf{Error} \\
\midrule
2024-03 & 2024-04 & 39 & 129{,}390 & 129{,}639 & \(-0{,}249\) \\
2024-04 & 2024-05 & 40 & 129{,}860 & 130{,}123 & \(-0{,}263\) \\
2024-05 & 2024-06 & 41 & 129{,}570 & 130{,}587 & \(-1{,}017\) \\
\multicolumn{6}{c}{\dots (21 orígenes intermedios) \dots} \\
2026-03 & 2026-04 & 63 & 144{,}310 & 144{,}534 & \(-0{,}224\) \\
2026-04 & 2026-05 & 64 & 144{,}590 & 145{,}001 & \(-0{,}411\) \\
\bottomrule
\end{tabular}
\end{table}

Las métricas del ejemplo se calculan sobre los 26 orígenes completos, no sobre
la muestra mostrada.

\section{Cálculo de métricas}

\subsection{Definiciones}

\begin{align}
\mathrm{MAE}&=\frac{1}{n}\sum_{i=1}^{n}|y_i-\hat y_i|, &
\mathrm{RMSE}&=\sqrt{\frac{1}{n}\sum_{i=1}^{n}(y_i-\hat y_i)^2},\\
\mathrm{MAPE}&=\frac{100}{n}\sum_{i=1}^{n}\left|\frac{y_i-\hat y_i}{y_i}\right|, &
\mathrm{sMAPE}&=\frac{100}{n}\sum_{i=1}^{n}\frac{|y_i-\hat y_i|}{(|y_i|+|\hat y_i|)/2}.
\end{align}

El MASE escala el error con la referencia \emph{naive} calculada dentro de la
ventana de entrenamiento de cada origen, de modo que la escala nunca incorpora
información del tramo evaluado \cite{hyndman_koehler2006}:
\begin{equation}
\mathrm{MASE}=\frac{1}{m}\sum_{r}\frac{|e_{r,h}|}{s_r},
\qquad
s_r=\frac{1}{r-1}\sum_{t=2}^{r}|y_t-y_{t-1}| .
\end{equation}

\subsection{Sustitución numérica}

Con los tres primeros errores de la tabla anterior
(\(-0{,}249\), \(-0{,}263\), \(-1{,}017\)), el aporte parcial al MAE es
\begin{equation}
\frac{|-0{,}249|+|-0{,}263|+|-1{,}017|}{3}=\frac{1{,}529}{3}=0{,}510 ,
\end{equation}
y al RMSE
\begin{equation}
\sqrt{\frac{0{,}249^{2}+0{,}263^{2}+1{,}017^{2}}{3}}
=\sqrt{\frac{0{,}062+0{,}069+1{,}034}{3}}=0{,}618 .
\end{equation}
Para el MAPE, el primer término es
\(|-0{,}249/129{,}390|\times100=0{,}193\%\).

Extendiendo el cálculo a los 26 orígenes, la aplicación obtiene los valores de
la Tabla~\ref{tab:metricas}. El MASE de \(1{,}010\) indica que el error del
modelo es prácticamente igual al de la referencia naive dentro de muestra; se
reporta como métrica auxiliar y no bloquea por sí solo, porque la comparación
decisiva se hace contra los benchmarks en backtesting.

\begin{table}[H]
\centering
\caption{Métricas de backtesting por modelo en \(h=1\) (26 orígenes), ordenadas por RMSE}
\label{tab:metricas}
\begin{tabular}{@{}lrrrrr@{}}
\toprule
\textbf{Modelo} & \textbf{MAE} & \textbf{RMSE} & \textbf{MAPE (\%)} & \textbf{sMAPE (\%)} & \textbf{MASE} \\
\midrule
\textbf{Drift} & \textbf{0{,}7058} & \textbf{0{,}9295} & \textbf{0{,}5221} & \textbf{0{,}5229} & \textbf{1{,}0099} \\
Log-variación & 0{,}7652 & 0{,}9896 & 0{,}5674 & 0{,}5681 & 1{,}0896 \\
Variación mensual & 0{,}7679 & 0{,}9910 & 0{,}5695 & 0{,}5702 & 1{,}0933 \\
Naive & 0{,}7012 & 1{,}0959 & 0{,}5128 & 0{,}5159 & 1{,}0045 \\
Holt lineal & 0{,}8770 & 1{,}1671 & 0{,}6462 & 0{,}6479 & 1{,}2454 \\
\bottomrule
\end{tabular}
\end{table}

Obsérvese que Naive tiene un MAE ligeramente menor que Drift
(\(0{,}7012\) frente a \(0{,}7058\)) pero un RMSE claramente mayor
(\(1{,}0959\) frente a \(0{,}9295\)): Naive comete errores grandes con más
frecuencia, porque no incorpora la tendencia. El RMSE penaliza esos errores
grandes de forma cuadrática, y es el criterio que usa la selección.

\section{Selección del modelo}

La aplicación elige \textbf{un único modelo por serie} para toda la
trayectoria. Cada modelo \(k\) se puntúa por su RMSE relativo al mejor de cada
horizonte, ponderando cada horizonte por \(1/h\):
\begin{equation}
S_k=\frac{\displaystyle\sum_{h}\frac{1}{h}\cdot\frac{\mathrm{RMSE}_{k,h}}{\min_j \mathrm{RMSE}_{j,h}}}
{\displaystyle\sum_{h}\frac{1}{h}},
\qquad
k^{*}=\arg\min_k S_k .
\end{equation}

La ponderación \(1/h\) refleja que el error crece con el horizonte y que la
evidencia de corto plazo es más confiable. La ventana de agregación son todos
los horizontes evaluados, que dependen únicamente de la longitud de la serie;
por eso el modelo ---y con él los primeros valores proyectados--- no cambia
según el horizonte que solicite el usuario.

\paragraph{Resultado.} Gana \textbf{Drift} (\(S_k\) mínimo), con el menor RMSE
en \(h=1\) y un comportamiento estable en horizontes mayores. Los modelos de
suavizamiento de Holt quedaron por detrás porque reaccionan al último
movimiento y esta serie tiene ruido mensual apreciable; las regresiones sobre
el nivel resultaron menos precisas en el tramo corto. El modelo seleccionado se
fija \emph{antes} de clasificar los horizontes, de manera que la admisibilidad,
las métricas, los intervalos y la trayectoria correspondan todos a Drift.

\section{Determinación del horizonte viable}

La aplicación evalúa horizontes mensuales consecutivos y clasifica cada uno.
Un horizonte solo es defendible si todos los anteriores lo son.

\begin{table}[H]
\centering
\caption{Resumen de la decisión de horizonte}
\begin{tabular}{@{}lr@{}}
\toprule
\textbf{Concepto} & \textbf{Valor} \\
\midrule
Horizonte solicitado & 18 meses \\
Horizontes evaluados & \(h=1\) a \(h=24\) \\
Máximo evaluable por datos & 24 meses \\
Máximo con evidencia completa (\(\geq 6\) ventanas) & 21 meses \\
\textbf{Máximo recomendado} & \textbf{12 meses} \\
\textbf{Máximo admisible (incluye escenario)} & \textbf{24 meses} \\
Primer horizonte no viable & Ninguno \\
Estado del horizonte solicitado & Escenario de alta incertidumbre \\
\bottomrule
\end{tabular}
\end{table}

\begin{table}[H]
\centering
\caption{Clasificación por horizonte (muestra)}
\begin{tabular}{@{}rlrrrl@{}}
\toprule
\textbf{h} & \textbf{Modelo} & \textbf{Ventanas} & \textbf{MAPE (\%)} & \textbf{IC95 rel.} & \textbf{Clasificación} \\
\midrule
1 & Drift & 26 & 0{,}522 & 2{,}28\% & Técnica con cautela \\
2 & Drift & 25 & 0{,}954 & 3{,}62\% & Técnica con cautela \\
10 & Drift & 17 & 1{,}264 & 4{,}40\% & Extendida con cautela \\
12 & Drift & 15 & 1{,}142 & 3{,}77\% & Extendida con cautela \\
13 & Drift & 14 & 1{,}518 & 4{,}16\% & Escenario \\
18 & Drift & 9 & 2{,}817 & 6{,}08\% & Escenario \\
\bottomrule
\end{tabular}
\end{table}

\paragraph{Por qué el corte está en 12 meses.} Los errores no se disparan: el
MAPE en \(h=18\) sigue siendo \(2{,}82\%\) y el IC95 relativo, \(6{,}08\%\),
muy por debajo de los umbrales de bloqueo. El corte se produce porque el modelo
seleccionado es un \emph{benchmark} (Drift) y, a partir de \(h=13\), la
aplicación degrada a escenario los horizontes largos sostenidos por un
benchmark simple: prolongar 18 meses una pendiente media histórica es
defendible como referencia, pero no como proyección técnica principal.

En consecuencia, la aplicación \textbf{no niega} el cálculo. Lo entrega con la
clasificación que corresponde: proyección técnica hasta 12 meses y escenario
entre 13 y 24.

\section{Cálculo de la proyección a 18 meses}

\subsection{Parámetros del modelo}

Drift extrapola la pendiente media entre el primer y el último valor:
\begin{equation}
\hat y_{T+h}=y_T+h\left(\frac{y_T-y_1}{T-1}\right).
\end{equation}

Con los valores reales de la serie:
\begin{align}
y_1&=100{,}801 \quad (\text{2021-01}), &
y_T&=144{,}590 \quad (\text{2026-05}), &
T&=65 .
\end{align}

La pendiente mensual es
\begin{equation}
b=\frac{y_T-y_1}{T-1}=\frac{144{,}590-100{,}801}{64}
=\frac{43{,}789}{64}=0{,}684203125\ \text{puntos de índice por mes}.
\end{equation}

\subsection{Sustitución para el horizonte solicitado}

\begin{equation}
\boxed{\;\hat y_{T+18}=144{,}590+18\times 0{,}684203125
=144{,}590+12{,}315656=156{,}905656\;}
\end{equation}

Comprobación intermedia para \(h=1\):
\(\hat y_{T+1}=144{,}590+0{,}684203=145{,}274203\), que coincide con el primer
valor de la trayectoria.

\subsection{Trayectoria completa}

\begin{longtable}{@{}rlrrrl@{}}
\caption{Trayectoria proyectada de \(h=1\) a \(h=18\) (modelo Drift)}\\
\toprule
\textbf{h} & \textbf{Periodo} & \textbf{Índice proyectado} & \textbf{IC95 inf.} & \textbf{IC95 sup.} & \textbf{Estado} \\
\midrule
\endfirsthead
\toprule
\textbf{h} & \textbf{Periodo} & \textbf{Proyectado} & \textbf{IC95 inf.} & \textbf{IC95 sup.} & \textbf{Estado} \\
\midrule
\endhead
1 & 2026-06 & 145{,}2742 & 144{,}1828 & 146{,}4314 & Técnica \\
2 & 2026-07 & 145{,}9584 & 144{,}4150 & 147{,}5949 & Técnica \\
3 & 2026-08 & 146{,}6426 & 144{,}7523 & 148{,}6469 & Técnica \\
4 & 2026-09 & 147{,}3268 & 145{,}1441 & 149{,}6412 & Técnica \\
5 & 2026-10 & 148{,}0110 & 145{,}5706 & 150{,}5985 & Técnica \\
6 & 2026-11 & 148{,}6952 & 146{,}0219 & 151{,}5297 & Técnica \\
7 & 2026-12 & 149{,}3794 & 146{,}4919 & 152{,}4410 & Extendida \\
8 & 2027-01 & 150{,}0636 & 146{,}9767 & 153{,}3366 & Extendida \\
9 & 2027-02 & 150{,}7478 & 147{,}4737 & 154{,}2194 & Extendida \\
10 & 2027-03 & 151{,}4320 & 147{,}9808 & 155{,}0913 & Extendida \\
11 & 2027-04 & 152{,}1162 & 148{,}4965 & 155{,}9542 & Extendida \\
12 & 2027-05 & 152{,}8004 & 149{,}0198 & 156{,}8090 & Extendida \\
13 & 2027-06 & 153{,}4846 & 149{,}5496 & 157{,}6569 & Escenario \\
14 & 2027-07 & 154{,}1688 & 150{,}0853 & 158{,}4986 & Escenario \\
15 & 2027-08 & 154{,}8530 & 150{,}6262 & 159{,}3348 & Escenario \\
16 & 2027-09 & 155{,}5373 & 151{,}1717 & 160{,}1660 & Escenario \\
17 & 2027-10 & 156{,}2215 & 151{,}7216 & 160{,}9926 & Escenario \\
\textbf{18} & \textbf{2027-11} & \textbf{156{,}9057} & \textbf{152{,}2753} & \textbf{161{,}8151} & \textbf{Escenario} \\
\bottomrule
\end{longtable}

Como Drift es lineal en \(h\), la diferencia entre meses consecutivos es
constante e igual a \(b=0{,}6842\).

\section{Resultado del índice proyectado}

\begin{table}[H]
\centering
\caption{Resultado entregado por la aplicación}
\begin{tabular}{@{}lr@{}}
\toprule
\textbf{Campo} & \textbf{Valor} \\
\midrule
Modelo aplicado & Drift \\
Periodo proyectado & 2027-11 \\
\textbf{Índice proyectado (\(h=18\))} & \textbf{156{,}9057} \\
Intervalo de predicción del 95\,\% & [152{,}2753;\ 161{,}8151] \\
Factor de actualización & 1{,}0852 \\
Variación acumulada proyectada & 8{,}518\% \\
Estado & Escenario de alta incertidumbre \\
Nivel de confianza & Bajo \\
\midrule
Último valor \textbf{recomendado} (\(h=12\)) & 152{,}8004 (2027-05) \\
\bottomrule
\end{tabular}
\end{table}

\paragraph{Respuesta directa.} El valor del índice proyectado para
\textit{Vías urbanas} a 18 meses es \textbf{156{,}9057}, correspondiente a
noviembre de 2027, con un intervalo de 95\% entre \(152{,}28\) y \(161{,}82\).

\paragraph{Por qué no debe usarse como valor recomendado.} Este valor está
clasificado como \emph{escenario}: procede de extender 18 meses la pendiente
media histórica mediante un benchmark. Si se requiere un valor con respaldo de
proyección técnica, debe usarse el horizonte de 12 meses
(\(152{,}8004\), mayo de 2027). El valor a 18 meses es una referencia
orientativa que debe leerse junto con su banda y acompañarse del criterio del
profesional responsable.

\section{Interpretación técnica}

La serie crece de forma sostenida (\(43{,}44\%\) acumulado en 65 meses) con
ruido mensual apreciable. En ese contexto:

\begin{itemize}
  \item Drift resulta preferible porque captura la tendencia de largo plazo sin sobrerreaccionar a los movimientos recientes: obtiene el menor RMSE en el tramo corto (\(0{,}9295\) frente a \(1{,}0959\) de Naive).
  \item El error crece de forma ordenada con el horizonte: el MAPE pasa de \(0{,}52\%\) en \(h=1\) a \(2{,}82\%\) en \(h=18\), y el IC95 relativo de \(2{,}28\%\) a \(6{,}08\%\). No hay señales de inestabilidad.
  \item El límite de 12 meses no proviene de un mal desempeño sino de una regla de prudencia: no se sostiene una proyección técnica de largo plazo sobre un benchmark simple.
  \item La proyección supone que las condiciones que generaron la tendencia se mantienen. Un cambio de régimen de precios invalidaría la extrapolación, y esa hipótesis no es verificable estadísticamente con la serie.
\end{itemize}

\section{Reproducibilidad del cálculo}

\begin{table}[H]
\centering
\caption{Elementos para reproducir el ejemplo}
\begin{tabular}{@{}p{4.6cm}p{10.4cm}@{}}
\toprule
\textbf{Elemento} & \textbf{Valor} \\
\midrule
Archivo de datos & \texttt{anex-ICOCIV-may2026.xlsb} (raíz del repositorio) \\
Tabla de origen & \texttt{T\_16\_2} (A16.2 Índice por tipologías de obra) \\
Ruta & Carreteras\dots\ $\rightarrow$ Carreteras (excepto elevadas); calles $\rightarrow$ Vías urbanas \\
Función de carga & \texttt{app\_icociv.datos.cargador\_datos.cargar\_todas\_tablas} \\
Resolución de fila & \texttt{app\_icociv.proyeccion.servicio\_proyeccion.resolver\_fila\_seleccionada} \\
Construcción de serie & \texttt{app\_icociv.proyeccion.servicio\_proyeccion.construir\_serie} \\
Proyección & \texttt{app\_icociv.proyeccion.servicio\_proyeccion.ejecutar\_proyeccion} \\
Backtesting & \texttt{app\_icociv.validacion.backtesting.ejecutar\_backtesting} \\
Parámetros de llamada & \texttt{ejecutar\_proyeccion(serie, 2027, 11, 2021)} \\
Archivo de salida & \texttt{docs/ejemplo\_vias\_urbanas\_proyeccion.csv} \\
Dependencias & pandas, NumPy, scikit-learn, pyxlsb, openpyxl, matplotlib, PySide6 \\
\bottomrule
\end{tabular}
\end{table}

El archivo CSV adjunto contiene las 65 observaciones históricas y los 18
valores proyectados con las columnas \texttt{periodo},
\texttt{valor\_observado}, \texttt{valor\_proyectado}, \texttt{horizonte},
\texttt{modelo} y \texttt{estado}.

\paragraph{Verificación manual mínima.} Basta con los dos extremos de la serie:
\begin{equation}
\hat y_{T+18}=144{,}590+18\times\frac{144{,}590-100{,}801}{64}=156{,}9057 .
\end{equation}

\section*{Referencias}
\addcontentsline{toc}{section}{Referencias}
\begin{thebibliography}{9}

\bibitem{dane_metodologia_icociv}
DANE. \textit{Metodología General Índice de Costos de la Construcción de Obras
Civiles --- ICOCIV}. Departamento Administrativo Nacional de Estadística.

\bibitem{hyndman2021fpp3}
Hyndman, R.~J. y Athanasopoulos, G. (2021). \textit{Forecasting: Principles and
Practice} (3.\textsuperscript{a} ed.). OTexts, Melbourne.

\bibitem{hyndman_koehler2006}
Hyndman, R.~J. y Koehler, A.~B. (2006). Another look at measures of forecast
accuracy. \textit{International Journal of Forecasting}, 22(4), 679--688.

\bibitem{tashman2000}
Tashman, L.~J. (2000). Out-of-sample tests of forecasting accuracy: an analysis
and review. \textit{International Journal of Forecasting}, 16(4), 437--450.

\bibitem{holt1957}
Holt, C.~C. (1957). \textit{Forecasting seasonals and trends by exponentially
weighted moving averages}. Carnegie Institute of Technology.

\bibitem{gardner_mckenzie1985}
Gardner, E.~S. y McKenzie, E. (1985). Forecasting trends in time series.
\textit{Management Science}, 31(10), 1237--1246.

\bibitem{granger_newbold1974}
Granger, C.~W.~J. y Newbold, P. (1974). Spurious regressions in econometrics.
\textit{Journal of Econometrics}, 2(2), 111--120.

\bibitem{nist_outliers}
NIST/SEMATECH. \textit{e-Handbook of Statistical Methods}: detección y manejo de
valores atípicos.

\end{thebibliography}

\end{document}
